A functor $H : C \rightarrow D$ is said to \textit{preserve the limits} of a functor $F : J \rightarrow C$ when every limiting cone $\nu : b \rightarrow F$ in $C$ for $F$ yields by composition with $H$ a limiting cone $H\nu : Hb \rightarrow HF$ in $D$. This requires not only that $H$ take each limit object which exists in $C$ to a limit object in $D$ but also that $H$ take limiting cones to limiting cones.
